\section{第一段创业：推荐系统、秒针与 Web2 技术复杂度}

本章承接导读里的“公司史也是技术商业案例”这一判断，先回答一个问题：为什么早期 Web2 数据能力会成为企业级 AI 的伏笔。吴明辉最开始的商业计划书是 recommendation system，而不是广告系统；这个细节很重要，因为它说明团队最早关心的是 Web2 时代个性化内容、用户 ID、推荐和动态系统。Web1 的页面可以缓存，所有人看同一份内容；Web2 以后，每个用户看到自己的内容，系统复杂度明显上升。

\begin{importantbox}{Web2 给技术团队带来的新问题}
当每个人看到的内容都不同，缓存、数据库、个性化、推荐、日志和实时计算都会变复杂。早期推荐系统和广告监测不是今天意义上的大模型，但它们已经在处理“数据驱动的互联网行为理解”。
\end{importantbox}

\begin{knowledgebox}{术语消化：从 Web2 到企业 AI 的连续性}
\begin{tabularx}{\textwidth}{p{0.22\textwidth}p{0.32\textwidth}X}
\toprule
概念 & 当年解决的问题 & 为什么和今天 AI 有关 \\
\midrule
User ID & 让系统知道“谁在看” & 企业 AI 也必须知道任务属于谁、权限属于谁、上下文属于谁。 \\
推荐系统 & 为不同用户分配不同内容 & Agentic Model 也要为不同角色生成不同动作建议。 \\
广告监测 & 让多方相信同一份数据 & 企业 AI 输出必须可审计，否则无法进入经营决策。 \\
日志与行为数据 & 记录真实互动 & 私有数据的价值来自持续记录业务过程。 \\
\bottomrule
\end{tabularx}
\end{knowledgebox}

从这个角度看，第一段创业不是“还没进入 AI 的前传”，而是企业级 AI 的底层训练。推荐系统训练的是个性化分发能力，广告监测训练的是可信计量能力，Web2 基础设施训练的是高并发和用户上下文能力。这些能力后来换了名字：个性化变成 context，可信计量变成 auditability，高并发系统变成企业 workflow 的稳定交付。

\begin{knowledgebox}{从旧能力到新能力的迁移}
\begin{tabularx}{\textwidth}{p{0.24\textwidth}p{0.32\textwidth}X}
\toprule
旧能力 & 当年的业务形态 & AI 时代的迁移 \\
\midrule
个性化推荐 & 给不同用户看不同内容 & 给不同角色、部门和客户生成不同任务建议。 \\
广告监测 & 衡量流量、曝光和转化 & 衡量 Agent 动作是否真实带来业务结果。 \\
中立第三方 & 多方认可同一套数据 & 企业 AI 输出要能被销售、法务、财务和客户共同接受。 \\
外包/交付训练 & 给新浪、搜狐等客户做系统 & 理解真实客户需求，而不是只做实验室 demo。 \\
\bottomrule
\end{tabularx}
\end{knowledgebox}

\begin{knowledgebox}{第一段创业的三条隐线}
\begin{tabularx}{\textwidth}{p{0.20\textwidth}p{0.34\textwidth}X}
\toprule
隐线 & 当时看起来是什么 & 后来看起来是什么 \\
\midrule
技术隐线 & 给门户、广告和客户系统做工程交付 & 学会在真实业务系统里处理数据和可靠性。 \\
商业隐线 & 从外包、推荐系统转向广告监测 & 学会在多方交易中找到可信中间层。 \\
心理隐线 & 数学竞赛和创业随机性 & 学会把长期不确定性当作可解问题。 \\
\bottomrule
\end{tabularx}
\end{knowledgebox}

\subsection{秒针的成功：从技术到中立测量}

秒针系统后来成功，是因为广告主、媒体和代理之间需要一个中立测量方。这个角色的价值在于数据可信、标准统一和商业中立。吴明辉也提到，如果秒针自己下场包流量，就会失去中立裁判身份。这是 ToB 公司很典型的边界问题：短期收入诱惑可能损害长期信任。

\begin{knowledgebox}{中立测量的商业逻辑}
广告监测公司的客户不是要一个“看起来更聪明”的算法，而是要一个可被多方承认的事实基础。可信数据一旦成为交易标准，就比单点功能更有壁垒。
\end{knowledgebox}

\subsection{从数学竞赛到创业心理}

访谈中多次提到数学竞赛训练、奥林匹克竞技和心理调适。它解释了吴明辉的创业耐力：把困难看成一道题，相信花更长时间也能解出来。这种思维对 ToB 创业很关键，因为 ToB 很少一夜爆发，更多是长周期交付、客户信任和组织韧性。

\begin{warningbox}{竞赛思维的双刃剑}
竞赛训练能带来强问题拆解能力和延迟满足能力，但也可能让创始人低估市场、销售、组织和资本结构的复杂性。ToB 公司不是只解技术题，还要解客户、现金流和团队题。
\end{warningbox}

\subsection{本章小结}

第一段创业说明，吴明辉不是从“AI 概念”进入企业服务，而是从 Web2 数据、广告监测、中立测量和推荐系统进入。这些早期能力后来转化为企业私有数据和 Agentic Model 的底层直觉。

